\subsection{The raw pentagon leaves six normalized candidates}

\begin{statement}{D23.P26}{Proved}{An exact six-candidate necessary reduction}
\end{statement}
Normalize a hypothetical actual full-GT element with nonzero
$\Delta$ by D23.H2 to $(\rho,a_3,a_2)=(0,1,0)$.
Use the zero-based indices in the array
\texttt{mixed72\_V\_1} of D23.P22. Its degree-$1344$
conjugating permutation must have one of the six indices
\[
                         1,19,23,27,53,61.
\]
For each such candidate exactly three of D23.H8's mixed
values survive the literal degree-$8064$ pentagon. The full
retained list, indexed in D23.H8's ordered $216$-element
centralizer subgroup, is
\[
\begin{array}{c|l}
V\text{ index}&\text{mixed centralizer indices}\\\hline
1&109,129,137\\
19&112,120,140\\
23&40,48,68\\
27&44,52,60\\
53&117,125,133\\
61&81,89,97
\end{array}
\]
Thus there are eighteen remaining pairs $(V,R)$ in this
necessary test. No assertion that any pair is realized by
an actual gentle or full-GT element is made.

\begin{proof}
We first justify why a full-GT element must satisfy the
\emph{literal} pentagon in the raw braid action, although
that action does not factor through $B_4/\langle z_4\rangle$.
Write $a,b,c$ for the positive braid generators and put
\[
\begin{aligned}
X_1&=a^2,&X_2&=b^2,&X_3&=c^2,\\
X_4&=(aba)^2,&X_5&=(cbc)^2,&z_4&=(abc)^4.
\end{aligned}
\]
Supplied FP.30.1 and FP.30.3 identify their images with
the five standard cyclic boundary twists and identify the
full-GT five-cycle in the order
\[
 \mathcal P(f)=f(X_3,X_4)f(X_5,X_1)f(X_2,X_3)
                    f(X_4,X_5)f(X_1,X_2).
                                                        \tag{D23.P26.1}
\]
The image of this product in
$\widehat{B_4/\langle z_4\rangle}$ is one.
The kernel of the natural map from $\widehat B_4$ is
the closed normal subgroup generated by $z_4$, hence its
procyclic central closure. It is canonically
$\widehat{\mathbb Z}$: the continuous exponent-sum map
$e:\widehat B_4\to\widehat{\mathbb Z}$ takes $z_4$ to
$12$, so the map $t\mapsto z_4^t$ is injective, as
multiplication by $12$ is injective on
$\widehat{\mathbb Z}=\prod_p\mathbb Z_p$. The same
central-kernel and exponent-sum argument is recorded in
supplied FP.26, paragraph \emph{Lifting the five-point
frame}.

Consequently $\mathcal P(f)=z_4^t$ for some
$t\in\widehat{\mathbb Z}$. Each factor has exponent
sum zero, since $f\in\widehat F_2'$ and an evaluation
followed by $e$ is an abelian homomorphism. Applying $e$
therefore gives $12t=0$, whence $t=0$. This proves
\[
                   \mathcal P(f)=1\quad\hbox{in }\widehat B_4.
                                                        \tag{D23.P26.2}
\]
In particular its degree-$8064$ raw image must be the
identity. This conclusion is stronger than identity
modulo the global deck group and requires no claim that
the raw action descends to the five-point quotient.

Next fix the complete first value. By D23.F36 an actual
normalized element with a candidate $V$ has
$F=f(X_1,X_2)=F_V$, uniquely, not merely modulo a
solvable quotient. The proof of F36 uses the trivial
center of the full perfect residual and the two-cycle,
so no unenumerated first-value fiber remains. D23.H7
and the exhaustive D23.H8 calculation then leave exactly
the twelve candidates and their three mixed values listed
in H8. These are all the possibilities after the necessary
base pentagon-pair membership test.

Let tildes denote the raw permutations. The projection of
the raw first-three-strand group to its degree-$1344$
image is a marked isomorphism by D23.P24, as is that of
$\langle\widetilde X_3,\widetilde X_4\rangle$.
For the remaining pair the direct stabilizer-chain
calculations give
\[
\begin{split}
|\langle\widetilde X_4,\widetilde X_5\rangle|
  &=|\langle X_4,X_5\rangle|=13374150672384,\\
|\langle\widetilde X_4,\widetilde X_5\rangle'|
  &=|\langle X_4,X_5\rangle'|=557256278016.
\end{split}                                             \tag{D23.P26.3}
\]
The projection is onto because it sends the two marked
generators to their base counterparts. Equality of finite
orders makes it an isomorphism, including on the derived
subgroups. Thus each possible base value $F,R,H$ has a
\emph{unique} marked raw lift. In particular an actual
evaluation of the same profinite $f$ must have exactly
those raw values; there is no separate scalar or deck
choice for an associator value.

Use $D=abcaba$, retaining its actual inverse. By D23.P25,
the base value of $H=f(X_4,X_5)$ is necessarily
$[D,F_VR]$, and H8 has checked its membership in the
full derived pair group. Write $\widetilde H$ for its
unique lift from that marked group. Equation (D23.P26.2)
is equivalent to the further literal identity
\[
 \widetilde H=
 [\widetilde D,\widetilde F_V\widetilde R].
                                                        \tag{D23.P26.4}
\]
The certificate computes both sides for every one of the
$12\cdot3=36$ complete H8 pairs. It evaluates the
strong-generator straight-line programs in the raw marked
generators, with every lifted permutation checked against
the base projection. It also directly multiplies the
original five factors in (D23.P26.1), so (D23.P26.4)
is not checked merely by a rearranged symbolic convention.

Exactly all three mixed values pass for each odd index
$1,19,23,27,53,61$. All three fail for each even partner
$0,18,22,26,52,60$. In every failing case the five-factor
product has order two, projects to the base identity,
and has fiber phases $0$ and $3$, both occurring.
It is therefore a nonconstant fiber permutation, not a
global power of the order-six deck transformation.
Both descriptions are recorded as additional exact checks;
nonidentity alone is enough to violate (D23.P26.2).
The identity associator gives the identity raw product as
an equal-sign control. The retained indices and mixed
values are precisely the displayed table.

The completeness assertion now follows from F36, H8,
and the unique raw lifts: any actual normalized mixed
element must appear among those thirty-six pairs, and
its raw pentagon removes exactly the rejected eighteen.
The eighteen survivors remain subject to all common-word
conditions involving the third pair and all other actual
gentle or full-GT conditions.
\end{proof}

\paragraph{Certificates and source locators.}
The pair-group orders and marked projection are rebuilt
by \path{agents/pentagon/hurwitz_raw_pair45.py}; its JSON
and NPZ retain both marked generator pairs. The complete
raw test is
\path{agents/pentagon/hurwitz_raw_pentagon45.py}, with
the JSON table and NPZ arrays of required values, unique
marked lifts, and all thirty-six literal pentagon products.
The strong-generator evaluator is
\path{agents/pentagon/hurwitz_marked_raw_lift.py}.
The base exhaustive fiber is
\path{code/hurwitz_pentagon45_fibers.py} and its records;
the first-value uniqueness input is
\path{agents/five_level_six/hurwitz_first_value_uniqueness.tex}.
The supplied universal marking and central-kernel input are
\path{sources/DELTA_ON_GT/agent_five/five_point_results.tex},
FP.30.1, FP.30.3, and FP.26's paragraph
\emph{Lifting the five-point frame}. The exact raw orbit
and its actual multiplier-one action are D23.H6.
